\documentclass[11pt,a4paper]{article}
\usepackage{gtu-el-yazisi}
\begin{document}
\notbilgi{Nümerik Analiz}{Ders 11 -- Lineer cebir denklem sistemleri için iterasyon yöntemleri}{numerik-11.hafta.pdf}

\sayfa{1}
\noindent\textit{ders 11}

\section{Lineer Cebir Denklem sistemleri için iterasyon yöntemleri}
\begin{enumerate}[label=\arabic*-]
  \item Matrisin özdeğer ve özvektörleri
  \item Jacobi iterasyon yöntemi
  \item Gauss-Seidel \quad '' \qquad ''
  \item iterasyon yöntemlerinin yakınsaklığı
\end{enumerate}

\subsection{Matris özdeğer ve özvektörleri, spektral yarıçap}
$n\times n$ boyutlu $A=[a_{ij}]$, \ $\lambda\in\mathbb{C}$ \ $\Rightarrow\ Ax=\lambda x$

$\Leftrightarrow\ (A-\lambda I)\,x=0$

$x=0 \rightarrow$ trivial (aşikar) çözümdür $\rightarrow$ bunu aramıyoruz

$x\neq 0 \rightarrow$ nontrivial \quad '' \quad $\checkmark \Rightarrow$ özdeğerler $\Rightarrow$ özvektörler

\medskip
$\det(A-\lambda I)\neq 0 \rightarrow$ trivial

$\det(A-\lambda I)=0 \rightarrow$ nontrivial $\checkmark$

$P(\lambda)=\det(A-\lambda I) \rightarrow$ $A$ matrisinin karakteristik polinomudur

$P(\lambda)$'nın sıfırları, özdeğerlerdir.
\[
  \underbrace{P(\lambda)=0}_{\lambda\ \text{özdeğer}}
  \qquad\Rightarrow\qquad
  \underbrace{(A-\lambda I)\,x=0}_{x\ \text{özvektörler}}
\]

\begin{ornek}
a)\quad $A=\begin{bmatrix}0&1\\-1&0\end{bmatrix} \qquad -\lambda\begin{bmatrix}1&0\\0&1\end{bmatrix}$
\[
  =\begin{bmatrix}-\lambda&1\\-1&-\lambda\end{bmatrix},\qquad \det(A-\lambda I)=\lambda^2+1
\]
\[ \lambda_{1,2}=\pm i \]
\[
  (A-\lambda_1 I)\,x=0 \ \Rightarrow\ \begin{bmatrix}-i&1\\-1&-i\end{bmatrix}\begin{bmatrix}x_1\\x_2\end{bmatrix}=0
\]
\[
  \left.\begin{aligned} -i\,x_1+x_2&=0\\ -x_1-i\,x_2&=0 \end{aligned}\right\}\ \kutula{x_2=i\,x_1}
  \ \Rightarrow\ \begin{bmatrix}x_1\\ i\,x_1\end{bmatrix} = x_1\begin{bmatrix}1\\ i\end{bmatrix}
\]
\end{ornek}

\sayfa{2}
$\lambda_2$ için \quad $x_1\begin{bmatrix}1\\-i\end{bmatrix}$
\[ \bar{X}_{1,2}=\begin{bmatrix}1\\ \pm i\end{bmatrix} \]

b)\quad $A=\begin{bmatrix}2&0&0\\1&1&2\\1&-1&4\end{bmatrix}$, \quad
$A-\lambda I=\begin{bmatrix}2&0&0\\1&1&2\\1&-1&4\end{bmatrix}-\lambda\begin{bmatrix}1&0&0\\0&1&0\\0&0&1\end{bmatrix}$
\begin{align*}
  P(\lambda)=\det(A-\lambda I)
  &=\begin{vmatrix}2-\lambda&0&0\\1&1-\lambda&2\\1&-1&4-\lambda\end{vmatrix}
   =(2-\lambda)\begin{vmatrix}1-\lambda&2\\-1&4-\lambda\end{vmatrix}\\
  &=\underbrace{(2-\lambda)}_{\lambda_1=2}\,\underbrace{(6-5\lambda+\lambda^2)}_{\lambda_{2,3}=2,3}=0
\end{align*}
\[ \lambda_{1,2}=2 \qquad \lambda_3=3 \]

\[
  (A-\lambda_{1,2}I)\,x=0 \qquad
  \begin{bmatrix}0&0&0\\1&-1&2\\1&-1&2\end{bmatrix}\begin{bmatrix}x_1\\x_2\\x_3\end{bmatrix}=\begin{bmatrix}0\\0\\0\end{bmatrix}
\]
\[ x_1-x_2+2x_3=0 \qquad \Rightarrow\ x_1=x_2-2x_3 \]
\[
  \begin{bmatrix}x_2-2x_3\\x_2\\x_3\end{bmatrix}
  =x_2\begin{bmatrix}1\\1\\0\end{bmatrix}+x_3\begin{bmatrix}-2\\0\\1\end{bmatrix}
\]
\[
  (A-\lambda_3 I)\,x=0 \ \Rightarrow\
  \begin{bmatrix}-1&0&0\\1&-2&2\\1&-1&1\end{bmatrix}\begin{bmatrix}x_1\\x_2\\x_3\end{bmatrix}=\begin{bmatrix}0\\0\\0\end{bmatrix}
\]
\[
  \begin{aligned}
    &\kutula{-x_1=0}\\
    &x_1-2x_2+2x_3=0 \qquad\Rightarrow\ x_2=x_3\\
    &x_1-x_2+x_3=0
  \end{aligned}
\]
\[ x_1=0,\quad x_2=x_3, \]

\sayfa{3}
\[
  \lambda_3=3 \ \Rightarrow\ \begin{bmatrix}x_1\\x_2\\x_3\end{bmatrix}=\begin{bmatrix}0\\x_3\\x_3\end{bmatrix}=x_3\begin{bmatrix}0\\1\\1\end{bmatrix}
\]

$A$'nın \quad $\lambda_1,\lambda_2,\dots,\lambda_n$ \quad özdeğerinin
\[ r(A)=\max\bigl\{\bigl(|\lambda_1|,|\lambda_2|,|\lambda_3|,\dots,|\lambda_n|\bigr)\bigr\} \]

1. örnekte $r(A)=1$\\
2. örnekte $r(A)=3$

\subsection{Yakınsak Matrisler}
$n\times n$ boyutlu \quad $A=[a_{ij}]$
\[ A^2=A\,A,\qquad A^3=A\cdot A\cdot A \]
\[
  A^k=\underbrace{A\cdot A\cdots A}_{k} \ \longrightarrow\ [A^k]_{ij} :\ A^k\text{'nın } i\text{'inci satırı, } j\text{'inci sütunu}
\]

Eğer her \ $i,j=1,\dots,n$ \quad $\displaystyle\lim_{k\to\infty}[A^k]_{ij}=0$ \
olursa $A$ matrisi yakınsak matristir.

\[ A=\begin{bmatrix}\frac12&0\\[1mm]\frac14&\frac12\end{bmatrix} \qquad \text{yakınsak old. gösteriniz} \]
\[
  A^2=A\,A=\begin{bmatrix}\frac12&0\\[1mm]\frac14&\frac12\end{bmatrix}\begin{bmatrix}\frac12&0\\[1mm]\frac14&\frac12\end{bmatrix}
  =\begin{bmatrix}\frac14&0\\[1mm]\frac14&\frac14\end{bmatrix}
\]
\[
  A^3=A\,A^2=\begin{bmatrix}\frac12&0\\[1mm]\frac14&\frac12\end{bmatrix}\begin{bmatrix}\frac14&0\\[1mm]\frac14&\frac14\end{bmatrix}
  =\begin{bmatrix}\frac18&0\\[1mm]\frac{3}{16}&\frac18\end{bmatrix}
\]
\[
  A^k=\begin{bmatrix}\frac{1}{2^k}&0\\[1mm]\frac{k}{2^{k+1}}&\frac{1}{2^k}\end{bmatrix}
\]
\[
  \lim_{k\to\infty}\frac{1}{2^k}=0 \quad\text{ve}\quad \lim_{k\to\infty}\frac{k}{2^{k+1}}=0
\]
\[ \Longrightarrow\ A \text{ yakınsak matristir} \]

\sayfa{4}
$A$ matrisinin yakınsak olması için, onun spektral yarıçapının $r(A)<1$ olması ile gerek ve yeter
koşul sağlanmış olur.

\subsection{2- Jacobi iterasyon yöntemi}
$\rightarrow$ $n\times n$ boyutlu $A=[a_{ij}]$ olsun öyle ki: \quad $a_{ii}\neq 0$ \quad
$x=[x_1,x_2,\dots,x_n]^T$ bilinmeyen ve $b=[b_1,b_2,\dots,b_n]^T$ sağlandığı için
\[
  Ax=b \ \Leftrightarrow\
  \begin{cases}
    a_{11}x_1+a_{12}x_2+\dots+a_{1n}x_n=b_1\\
    a_{21}x_1+a_{22}x_2+\dots+a_{2n}x_n=b_2\\
    \quad\vdots\\
    a_{n1}x_1+a_{n2}x_2+\dots+a_{nn}x_n=b_n
  \end{cases}
\]
\begin{align*}
  x_1&=-\frac{a_{12}}{a_{11}}x_2-\dots-\frac{a_{1n}}{a_{11}}x_n+\frac{b_1}{a_{11}}\\
  x_2&=-\frac{a_{21}}{a_{22}}x_1-\dots-\frac{a_{2n}}{a_{22}}x_n+\frac{b_2}{a_{22}}\\
  &\ \ \vdots\\
  x_n&=-\frac{a_{n1}}{a_{nn}}x_1-\dots-\frac{a_{n,n-1}}{a_{nn}}x_n\footnotemark+\frac{b_n}{a_{nn}}
\end{align*}
\footnotetext{Claude notu: muhtemelen $x_{n-1}$ olmalı.}
\[
  x_i=\sum_{\substack{j=1\\ j\neq i}}^{n}\left(-\frac{a_{ij}}{a_{ii}}\,x_j\right)+\frac{b_i}{a_{ii}}
  \qquad i=1,2,\dots,n
\]
denklem sistemi

\medskip
denklemlerin çözümü için yaklaşımı (iteratif yöntem)
\[
  x^{(0)}=0 \quad\Rightarrow\quad x^{(0)}=\begin{bmatrix}0\\0\\0\end{bmatrix}
\]
\[
  x_i^{(k)}=\frac{1}{a_{ii}}\left[-\sum_{\substack{j=1\\ j\neq i}}^{n}a_{ij}\,x_j^{(k-1)}+b_i\right]
\]
\[ i=1,2,\dots,n \]

\sayfa{5}
\begin{ornek}[3]
\[
  \begin{aligned}
    3x_1-x_2+x_3&=1\\
    3x_1+6x_2+2x_3&=0\\
    3x_1+3x_2+7x_3&=0
  \end{aligned}
\]
denklem sistemi için $x^{(0)}=0$ olacak biçimde Jacobi iterasyonu uygulayıp 2 iterasyonu bulunuz.
\end{ornek}
\[
  \left.\begin{aligned}
    x_1-\tfrac13x_2+\tfrac13x_3&=\tfrac13\\[1mm]
    \tfrac36x_1+x_2+\tfrac26x_3&=0\\[1mm]
    \tfrac37x_1+\tfrac37x_2+x_3&=0
  \end{aligned}\right\}
  \qquad
  \begin{aligned}
    x_1^{(k)}&=\tfrac13x_2^{(k-1)}-\tfrac13x_3^{(k-1)}+\tfrac13\\[1mm]
    x_2^{(k)}&=-\tfrac12x_1^{(k-1)}-\tfrac13x_3^{(k-1)}\\[1mm]
    x_3^{(k)}&=-\tfrac37x_1^{(k-1)}-\tfrac37x_2^{(k-1)}
  \end{aligned}
\]
\[ x^{(0)}=[0,\,0,\,0]^T \]
\[
  \Rightarrow\quad
  \begin{aligned}
    x_1^{(1)}&=\tfrac13\\ x_2^{(1)}&=0\\ x_3^{(1)}&=0
  \end{aligned}
  \qquad\qquad
  \Rightarrow\quad
  \begin{aligned}
    x_1^{(2)}&=\tfrac13\\ x_2^{(2)}&=-\tfrac16\\ x_3^{(2)}&=-\tfrac17
  \end{aligned}
\]

\sayfa{6}
\[
  \begin{bmatrix}a_{11}&a_{12}&\cdots&a_{1n}\\a_{21}&a_{22}&\cdots&a_{2n}\\ \vdots&&&\vdots\\a_{n1}&a_{n2}&\cdots&a_{nn}\end{bmatrix}
  =\underbrace{\begin{bmatrix}0&0&0&\cdots&0\\a_{21}&0&0&\cdots&0\\a_{31}&a_{32}&0&\cdots&0\\ \vdots&&&&\\a_{n1}&a_{n2}&\cdots&\cdots&0\end{bmatrix}}_{L}
  +\begin{bmatrix}a_{11}&0&0&0\\0&a_{22}&0&0\\0&0&\ddots&\\0&0&&a_{nn}\end{bmatrix}
\]
\[
  +\begin{bmatrix}0&a_{12}&\cdots&\cdots&a_{1n}\\0&0&a_{23}&\cdots&a_{2n}\\0&0&\cdots&\cdots&0\end{bmatrix}
\]
\begin{align*}
  \longrightarrow\quad A&=L+D+U\\
  Ax=b\ \Leftrightarrow\ (L+D+U)\,x&=b\\
  Dx&=-(L+U)\,x+b\\
  x&=\underbrace{-D^{-1}(L+U)}\,x+D^{-1}b
\end{align*}
\[ \Downarrow \]
\begin{align*}
  x&=T_J\,x+c_J\\
  T_J&=-D^{-1}(L+U)\\
  c_J&=D^{-1}b\\
  x^{(k)}&=T_J\,x^{(k-1)}+c_J
\end{align*}

\[
  \begin{bmatrix}3&-1&1\\3&6&2\\3&3&7\end{bmatrix}
  =\begin{bmatrix}3&0&0\\0&6&0\\0&0&7\end{bmatrix}+\begin{bmatrix}0&-1&1\\3&0&2\\3&3&0\end{bmatrix}
\]
\begin{align*}
  T_J&=-\begin{bmatrix}\frac13&0&0\\[1mm]0&\frac16&0\\[1mm]0&0&\frac17\end{bmatrix}\begin{bmatrix}0&-1&1\\3&0&2\\3&3&0\end{bmatrix}\\
  &=\begin{bmatrix}0&\frac13&-\frac13\\[1mm]-\frac12&0&-\frac13\\[1mm]-\frac37&-\frac37&0\end{bmatrix}
  \qquad c_J=\begin{bmatrix}\frac13\\[1mm]0\\0\end{bmatrix}
\end{align*}

\sayfa{7}
\subsection{3. Gauss-Seidel iterasyon yöntemi}
$\rightarrow$ $n\times n$ boyutlu $A=[a_{ij}]$ olsun öyle ki $a_{ii}\neq 0$\\
$x=[x_1,x_2,\dots,x_n]^T$ bilinmeyeni ve $b=[b_1,b_2,\dots,b_n]^T$ sağ tarafı için
\[
  Ax=b \ \Leftrightarrow\
  \begin{cases}
    a_{11}x_1+a_{12}x_2+\dots+a_{1n}x_n=b_1\\
    a_{21}x_1+a_{22}x_2+\dots+a_{2n}x_n=b_2\\
    \quad\vdots\\
    a_{n1}x_1+a_{n2}x_2+\dots+a_{nn}x_n=b_n
  \end{cases}
\]
\[ a_{i1}x_1+a_{i2}x_2+\dots+a_{in}x_n=b_i \]
\[
  x_i=\frac{1}{a_{ii}}\Biggl[-\underbrace{\sum_{j=1}^{i-1}a_{ij}\,x_j}_{\text{\okunamadi{k öğene}}}
  -\sum_{j=i+1}^{n}a_{ij}\,x_j+b_i\Biggr]
\]
\[
  x_i^{(k)}=\frac{1}{a_{ii}}\Biggl[-\sum_{j=1}^{i-1}a_{ij}\,x_j^{(k)}-\sum_{j=i+1}^{n}a_{ij}\,x_j^{(k-1)}+b_i\Biggr]
\]
\[ i=1,2,\dots,n \]

Gauss-Seidel \okunamadi{formülü}
\begin{align*}
  (L+D)\,x&=-Ux+b\\
  x&=(L+D)^{-1}Ux\footnotemark+(L+D)^{-1}b\\
  \det(L+D)&\neq 0
\end{align*}
\footnotetext{Claude notu: muhtemelen $x=-(L+D)^{-1}Ux+(L+D)^{-1}b$ olmalı (aşağıdaki $T_g$ ile uyumlu).}
\begin{align*}
  T_g&=-(L+D)^{-1}U\\
  c_g&=(L+D)^{-1}b
\end{align*}
\[ x=T_g\,x+c_g \]
\[ x^{(0)}=0 \]
\[ x^{(k)}=T_g\,x^{(k-1)}+c_g \]
\[ k=1,2,\dots \]

\sayfa{8}
\begin{ornek}
\noindent\begin{minipage}[c]{0.4\linewidth}
\[
  \begin{aligned}
    3x_1-x_2+x_3&=1\\
    3x_1+6x_2+4x_3&=0\\
    3x_1+3x_2+7x_3&=4
  \end{aligned}
\]
\end{minipage}\hfill
\begin{minipage}[c]{0.56\linewidth}
denklem sistemini $x^{(0)}=0$ olacak biçimde

Gauss-Seidel yöntemi uygulayarak ilk iki iterasyonunu bulunuz.
\end{minipage}
\end{ornek}

\noindent\begin{minipage}[c]{0.27\linewidth}
\small köşegen elemanını yalnız bırakıp normalize ederiz (genlik $=1$)
\end{minipage}\hfill
\begin{minipage}[c]{0.7\linewidth}
\[
  \left\{\begin{aligned}
    x_1&=\tfrac13x_2-\tfrac13x_3+\tfrac13\\[1mm]
    x_2&=-\tfrac12x_1-\tfrac23x_3\\[1mm]
    x_3&=-\tfrac37x_1-\tfrac37x_2+\tfrac47
  \end{aligned}\right.
\]
\end{minipage}

\medskip
\noindent\begin{minipage}[c]{0.27\linewidth}
\small köşegene kadar $k$, köşegenden sonra $(k-1)$ iterasyonları belirledik
\end{minipage}\hfill
\begin{minipage}[c]{0.7\linewidth}
\[
  \left\{\begin{aligned}
    x_1^{(k)}&=\tfrac13x_2^{(k-1)}-\tfrac13x_3^{(k-1)}+\tfrac13\\[1mm]
    x_2^{(k)}&=-\tfrac12x_1^{(k)}-\tfrac23x_3^{(k-1)}\\[1mm]
    x_3^{(k)}&=-\tfrac37x_1^{(k)}-\tfrac37x_2^{(k)}+\tfrac47
  \end{aligned}\right.
\]
\end{minipage}

\begin{align*}
  x_1^{(1)}&=\tfrac13\\[1mm]
  x_2^{(1)}&=-\tfrac12x_1^{(1)}=-\tfrac16\\[1mm]
  x_3^{(1)}&=-\tfrac37x_1^{(1)}-\tfrac37x_2^{(1)}+\tfrac47=\tfrac12
\end{align*}

$k=2$ olursa
\begin{align*}
  x_1^{(2)}&=-\tfrac13\left(-\tfrac16\right)-\tfrac13\left(\tfrac12\right)=+\tfrac13\footnotemark\\[1mm]
  x_2^{(2)}&=-\tfrac12x_1^{(2)}-\tfrac23\left(\tfrac12\right)\\[1mm]
  x_3^{(2)}&=-\tfrac37x_1^{(2)}-\tfrac37x_2^{(2)}+\tfrac47=\tfrac47
\end{align*}
\footnotetext{Claude notu: yukarıdaki formüle göre muhtemelen
$x_1^{(2)}=\tfrac13\left(-\tfrac16\right)-\tfrac13\left(\tfrac12\right)+\tfrac13=\tfrac19$ olmalı.}

\sayfa{9}
\[
  x^{(k)}=T_g\,x^{(k-1)}+c_g \qquad\qquad
  \begin{aligned} T_g&=-(L+D)^{-1}U\\ c_g&=(L+D)^{-1}b \end{aligned}
\]
\begin{align*}
  T_g:\quad x_1&=\tfrac13x_2-\tfrac13x_3+\tfrac13\\[1mm]
  x_2&=-\tfrac12\left(\tfrac13x_2-\tfrac13x_3+\tfrac13\right)-\tfrac23x_3
      =-\tfrac16x_2\ \okunamadi{$-\tfrac36x_3-\tfrac16$}\\[1mm]
  x_3&=-\tfrac37x_1-\tfrac37x_2+\tfrac47\\[1mm]
  &=-\tfrac1{14}x_2+\tfrac5{14}x_3+\tfrac12
\end{align*}
\[
  \begin{bmatrix}x_1\\x_2\\x_3\end{bmatrix}
  =\underbrace{\begin{bmatrix}0&\frac13&-\frac13\\[1mm]0&-\frac16&-\frac12\\[1mm]0&-\frac1{14}&\frac5{14}\end{bmatrix}}_{T_g\ \text{matrisi}}
  \begin{bmatrix}x_1\\x_2\\x_3\end{bmatrix}
  +\begin{bmatrix}\frac13\\[1mm]-\frac16\\[1mm]\frac12\end{bmatrix}
\]
\[ x^{(k)}=T\,x^{(k-1)}+c \qquad k=1,2 \]

\begin{teorem}
Her $x^{(0)}$ başlangıç yaklaşımı için $x^{(k)}=T\,x^{(k-1)}+c$, \ $k=1,2,\dots$ iterasyonuyla
belirlenmiş $x^{(k)}$, $k=0,1,\dots$ dizisinin $x=Tx+c$ denkleminin çözümüne yakınsaması için
gerek ve yeter koşul $r(T)<1$ eşitsizliğinin sağlanmalıdır.
\end{teorem}

\noindent\begin{minipage}[c]{0.5\linewidth}
\[
  \begin{aligned}
    2x_1-x_2+x_3&=-1\\
    2x_1+2x_2+2x_3&=4\\
    -x_1-x_2+2x_3&=-5
  \end{aligned}
\]
\end{minipage}\hfill
\begin{minipage}[c]{0.46\linewidth}
$T_J$, $T_g$ belirleyin ve yakınsaklığını inceleyin
\end{minipage}
\[
  T_J=\begin{bmatrix}0&\frac12&-\frac12\\[1mm]-1&0&-1\\[1mm]\frac12&\frac12&0\end{bmatrix}
  \qquad
  \left.\begin{aligned}
    \lambda_1&=0\\
    \lambda_2&=\tfrac{\sqrt5}{2}\,i\\
    \lambda_3&=-\tfrac{\sqrt5}{2}\,i
  \end{aligned}\right\}\ r(T_J)=\frac{\sqrt5}{2}>1
\]
\[
  T_g=\begin{bmatrix}0&\frac12&-\frac12\\[1mm]0&-\frac12&-\frac12\\ \phantom{0}&&\end{bmatrix}\footnotemark
\]
\footnotetext{Claude notu: notta $T_g$ matrisi yarım bırakılmış (son satır yazılmamış).}

\end{document}
